\problemname{Algebraic Teamwork}
The great pioneers of group theory and linear algebra want to cooperate and join their theories.
In group theory, permutations -- also known as bijective functions -- play an important role.
For a finite set $A$, a function $\sigma : A \rightarrow A$ is called a permutation of $A$ if and only if there is some function $
\rho : A → A$ with
$$\sigma(\rho(a)) = a\text{ and }\rho(\sigma(a)) = a\text{~~~for all }a \in A\text{.}$$
The other half of the new team -- the experts on linear algebra -- deal a lot with idempotent functions.
They appear as projections when computing shadows in 3D games or as closure operators like the transitive closure, just to name a few examples.
A function $p : A \rightarrow A$ is called idempotent if and only if
$$p(p(a)) = p(a)\text{~~~for all }a \in A\text{.}$$
To continue with their joined research, they need your help.
The team is interested in non-idempotent permutations of a given finite set $A$.
As a first step, they discovered that the result only depends on the set's size.
For a concrete size $1 \leq n \leq 10^5$, they want you to compute the number of permutations on a set of cardinality n that are not idempotent.
\begin{Input}
    The input starts with the number $t \leq 100$ of test cases.
    Then $t$ lines follow, each containing the set's size $1 \leq n \leq 10^5$.
\end{Input}

\begin{Output}
    Output one line for every test case containing the number modulo $1\,000\,000\,007 = (10^9 + 7)$ of \textbf{non}-idempotent permutations on a set of cardinality $n$.
\end{Output}
